% ==========================================================
%  EXAM QUESTIONS — DEEP REINFORCEMENT LEARNING
%  Shares the visual language of summary.tex (same accent,
%  same box style, same typography) but uses article class.
%
%  Coverage: credit assignment + the data-centric, planning,
%  model-based, and foundation-model material (lectures 06–11).
%  Real questions (IRL exam, credit-assignment quiz) are kept
%  verbatim; the rest are mock questions in the same style.
% ==========================================================

\documentclass[11pt,oneside]{article}

\usepackage[
  a4paper,
  top=2.4cm,
  bottom=2.1cm,
  left=2.3cm,
  right=2.3cm,
  headheight=15pt
]{geometry}

\usepackage[T1]{fontenc}
\usepackage[utf8]{inputenc}
\usepackage{lmodern}
\usepackage{microtype}
\usepackage{inconsolata}
\usepackage{amsmath,amssymb,mathtools,bm}
\usepackage{bbm}
\usepackage{enumitem}
\usepackage{booktabs,array,multicol,longtable}
\usepackage{xcolor}
\usepackage[most]{tcolorbox}
\usepackage{titlesec}
\usepackage{fancyhdr}
\usepackage[hidelinks]{hyperref}

\setlength{\parindent}{0pt}
\setlength{\parskip}{0.55em}
\setlist[itemize]{topsep=3pt,itemsep=2pt,leftmargin=2.2em}
\setlist[enumerate]{topsep=3pt,itemsep=2pt,leftmargin=2.2em}

% ── Colors ───────────────────────────────────────────────
\definecolor{accent}{HTML}{1B3A5C}
\definecolor{q@frame}{HTML}{1A5276}
\definecolor{q@bg}{HTML}{F6FAFD}
\definecolor{key@frame}{HTML}{1E8449}
\definecolor{key@bg}{HTML}{F4FBF6}
\definecolor{note@frame}{HTML}{9A6A00}
\definecolor{note@bg}{HTML}{FEF9E7}

% ── Box global style ─────────────────────────────────────
\tcbset{
  enhanced,
  breakable,
  arc=2.5mm,
  boxrule=0.5pt,
  left=4.5mm,
  right=4mm,
  top=2.5mm,
  bottom=2.5mm,
  fonttitle=\bfseries\small,
  sharp corners=west,
  parbox=false
}

\newtcolorbox{qbox}[1][]{
  colback=q@bg,
  colframe=q@frame,
  borderline west={3.5pt}{0pt}{q@frame},
  title={#1}
}

\newtcolorbox{keybox}[1][]{
  colback=key@bg,
  colframe=key@frame,
  borderline west={3.5pt}{0pt}{key@frame},
  title={#1}
}

\newtcolorbox{notebox}[1][]{
  colback=note@bg,
  colframe=note@frame,
  borderline west={3.5pt}{0pt}{note@frame},
  title={#1}
}

% ── Section headings ─────────────────────────────────────
\titleformat{\section}
  {\normalfont\Large\bfseries\color{accent}}
  {}{0em}{}
  [{\vspace{1pt}\color{accent!45}\titlerule[0.45pt]}]
\titlespacing*{\section}{0pt}{3.2ex plus 1ex minus .2ex}{1.8ex plus .2ex}

\titleformat{\subsection}
  {\normalfont\normalsize\bfseries\color{accent!80!black}}
  {}{0em}{}
\titlespacing*{\subsection}{0pt}{2.2ex plus 1ex minus .2ex}{1.0ex plus .2ex}

% ── Header & footer ──────────────────────────────────────
\pagestyle{fancy}
\fancyhf{}
\fancyhead[L]{\small\color{gray!70}Exam Questions — \coursetitle}
\fancyhead[R]{\small\color{gray!70}\thepage}
\renewcommand{\headrulewidth}{0.4pt}

% ── Document metadata ─────────────────────────────────────
\newcommand{\coursetitle}{Deep Reinforcement Learning}

% ── Question counter & macros ────────────────────────────
\newcounter{question}
\newcommand{\questiontitle}[1]{%
  \stepcounter{question}\begin{qbox}[Question \thequestion: #1]}
\newcommand{\closequestion}{\end{qbox}}

\newcommand{\repeatnote}{\textit{\color{gray!70}Asked multiple times.}}
\newcommand{\choice}{\item[$\square$] }
\newcommand{\answerline}[1]{\makebox[#1]{\hrulefill}}
\newcommand{\workingspace}[1]{%
  \par\vspace{2pt}%
  \begin{tcolorbox}[
    enhanced,
    breakable=false,
    colback=white,
    colframe=gray!40,
    boxrule=0.45pt,
    arc=0pt,
    sharp corners,
    height=#1,
    left=3mm, right=3mm, top=3mm, bottom=3mm,
    frame style={dashed}
  ]%
  \end{tcolorbox}%
  \vspace{2pt}%
}

% ── Math macros ──────────────────────────────────────────
\newcommand{\R}{\mathbb{R}}
\newcommand{\E}{\mathbb{E}}
\newcommand{\Prob}{\mathbb{P}}
\newcommand{\Var}{\mathrm{Var}}
\newcommand{\KL}[2]{D_{\mathrm{KL}}\!\left(#1\,\|\,#2\right)}
\newcommand{\dd}{\,\mathrm{d}}
\newcommand{\argmax}{\operatorname*{arg\,max}}
\newcommand{\argmin}{\operatorname*{arg\,min}}
\newcommand{\grad}{\nabla}
\newcommand{\T}{^{\mathsf{T}}}
\newcommand{\norm}[1]{\left\lVert #1 \right\rVert}
\newcommand{\abs}[1]{\left\lvert #1 \right\rvert}
\newcommand{\ind}{\mathbbm{1}}
\newcommand{\dataset}{\mathcal{D}}
\newcommand{\params}{\boldsymbol{\theta}}
\newcommand{\loss}{\mathcal{L}}
\newcommand{\Sset}{\mathcal{S}}
\newcommand{\Aset}{\mathcal{A}}
\newcommand{\Normal}{\mathcal{N}}
\newcommand{\bigO}[1]{\mathcal{O}\!\left(#1\right)}


% ==========================================================
%  BEGIN DOCUMENT
% ==========================================================
\begin{document}

% ── Title page ───────────────────────────────────────────
\thispagestyle{empty}
\vspace*{1.4cm}
\begin{center}
  {\color{accent}\rule{\linewidth}{1.8pt}}\par
  \vspace{0.65cm}
  {\fontsize{26}{32}\selectfont\bfseries\color{accent}Exam Questions\\[2pt] \coursetitle\par}
  \vspace{0.55cm}
  {\color{accent!35}\rule{\linewidth}{0.5pt}}\par
  \vspace{0.55cm}
  {\large Stefan Eder\par}
  \vspace{0.15cm}
  {\small\color{gray!70}\today\par}
\end{center}

\vspace{0.9cm}
\begin{notebox}[How this file is organized]
A practice set for the second half of the course: credit assignment, imitation learning
and offline RL, inverse RL, preference-based RL / RLHF, planning, model-based RL, and RL
for foundation models. \textbf{Real} old-exam questions (the Inverse RL exam and the
credit-assignment quiz) are reproduced verbatim with their answer choices intact; the
remaining \textbf{mock} questions follow the same style to broaden coverage. All answers
are in the \emph{Answer Key} at the end --- the question boxes are kept blank so the set
can be used for self-testing.
\end{notebox}

\tableofcontents
\newpage


% ==========================================================
\section{Credit Assignment (real quiz)}
% ==========================================================
\textit{\small From the credit-assignment / RUDDER quiz (DCA 2025), reproduced verbatim.}

\questiontitle{Reward transformations preserving optimality}
Which kind of transformation is \emph{guaranteed} to leave the set of optimal policies
unchanged?
\begin{itemize}[leftmargin=1.8em]
  \choice Clip to 1 all the rewards
  \choice Any affine transformation of immediate rewards only
  \choice A strictly monotonic transformation of the total return
  \choice Squaring the reward at each step
\end{itemize}
\closequestion

\questiontitle{Potential-based reward shaping}
The slides note that \emph{potential-based reward shaping} is mathematically equivalent to\dots
\begin{itemize}[leftmargin=1.8em]
  \choice doubling the learning rate
  \choice value-function normalization
  \choice reward clipping
  \choice initializing the $Q$-values with the potential function
\end{itemize}
\closequestion

\questiontitle{LSTM vs.\ FNN prediction errors}
Compared with a feed-forward network (FNN), an LSTM used for contribution analysis tends to
produce \emph{lower-variance} differences in consecutive predictions. What statistical
property of the LSTM's prediction errors explains this?
\begin{itemize}[leftmargin=1.8em]
  \choice Errors are \emph{highly correlated}, so they cancel in the difference
  \choice Errors have heavier tails, making large differences rare
  \choice Errors are i.i.d.\ across time
  \choice Errors are biased toward earlier time-steps
\end{itemize}
\closequestion

\questiontitle{Risk of online potential learning}
The slides note that potential-based shaping is \emph{equivalent to $Q$-value initialisation}
$Q'(s,a) = Q(s,a) + \Phi(s)$. What practical risk does this equivalence highlight when $\Phi$
is \emph{learned online} with TD-errors derived from sparse, delayed rewards?
\begin{itemize}[leftmargin=1.8em]
  \choice Bias is re-introduced because $\Phi$ itself must be bootstrapped from delayed feedback
  \choice The learning rate must be halved to keep $Q$-updates unbiased
  \choice Shaping may over-count rewards, double-updating the same transition
  \choice Eligibility traces converge to zero faster, causing under-estimation
\end{itemize}
\closequestion

\questiontitle{The credit-assignment problem (Sutton \& Minsky)}
According to R.\ Sutton and Minsky, what is the problem of credit assignment? (Select one or more.)
\begin{itemize}[leftmargin=1.8em]
  \choice Assigning credit or blame for overall episodes that contributed to the expectation of the return
  \choice Identifying models that get high rewards
  \choice Assigning credit or blame for overall outcomes to each of the decisions that contributed to those outcomes
  \choice Identifying actions that led to high rewards
\end{itemize}
\closequestion

\questiontitle{MC, TD, and reward delay}
Which of the following sentences are true? (Select one or more.)
\begin{itemize}[leftmargin=1.8em]
  \choice Variance in MC increases with the delay of the reward
  \choice TD methods correct the bias faster than MC
  \choice TD($\lambda$) assigns credit faster than TD($0$)
  \choice Variance in TD($0$) increases with the delay of the reward
\end{itemize}
\closequestion

\questiontitle{Zero future reward and the Markov property}
If we make the expected future rewards equal to zero (i.e.\ via reward redistribution), the
$Q$-values are just the average of the immediate reward. Why is this \emph{not} possible in a
standard MDP setting? (Select one or more.)
\begin{itemize}[leftmargin=1.8em]
  \choice Because to have future expected reward zero, states must be enriched with the past, which breaks the Markov assumption
  \choice It is completely possible in an MDP setting
  \choice Because in this case the reward depends also on the previous state and action, and is therefore not Markov
  \choice Because the reward is obtained before we enter the environment, and is therefore not causal, breaking the Markov property
\end{itemize}
\closequestion

\questiontitle{Definition of a sequence-Markov decision process (SDP)}
What is the definition of a sequence-Markov decision process (SDP), as seen in class?
(Select one or more.)
\begin{itemize}[leftmargin=1.8em]
  \choice A decision process equipped with a \textbf{Markov reward} and a \textbf{Markov policy}, but \textbf{transition probabilities} not required to be Markov
  \choice A decision process equipped with a \textbf{Markov policy}, but \textbf{transition probabilities and reward} not required to be Markov
  \choice A decision process equipped with a \textbf{Markov policy} and \textbf{Markov transition probabilities}, but a \textbf{reward} not required to be Markov
  \choice A decision process equipped with a \textbf{Markov reward} and \textbf{Markov transition probabilities}, but a \textbf{policy} not required to be Markov
\end{itemize}
\closequestion

\questiontitle{Return-equivalent SDPs}
What is the definition of \textbf{return-equivalent} SDPs? (Select one or more.)
\begin{itemize}[leftmargin=1.8em]
  \choice Two SDPs are return-equivalent if they differ only in their reward distributions but have the \textbf{same return per episode}
  \choice Two SDPs are return-equivalent if they differ only in their reward distributions but have the \textbf{same expected return}
  \choice Two SDPs are return-equivalent if they differ only in their reward distributions but have the \textbf{same expected return per episode}
\end{itemize}
\closequestion

\questiontitle{The explaining-away problem}
What is the \textbf{explaining-away problem} in MDPs, when a function $g$ is used to predict
the return based on the sequence of actions and states? (Select one or more.)
\begin{itemize}[leftmargin=1.8em]
  \choice Since later states are used for return prediction, earlier actions causing reward are missed
  \choice Previous actions can explain away future states, breaking the Markov property
  \choice Since the sequence of states and actions are mutually independent, the prediction needs every component to explain the return
  \choice In an MDP, one only needs the difference of states/actions to make a prediction, not the whole sequence
\end{itemize}
\closequestion

\questiontitle{LSTM vs.\ FFN for prediction}
Which statements are true regarding differences between LSTM and feed-forward networks (FFN)
for making future predictions? (Select one or more.)
\begin{itemize}[leftmargin=1.8em]
  \choice An LSTM will only learn to store information if a state-action pair has strong evidence for a change in the expected return
  \choice FFNs are not suited for processing sequences
  \choice LSTMs tend to have smaller prediction errors since they reuse past information
  \choice FFNs are prone to prediction errors since they must extract information from every state to make a prediction at every timestep
  \choice Prediction errors of LSTMs are more likely to cancel via prediction differences, since consecutive predictions rely on the same internal state and are highly correlated
\end{itemize}
\closequestion

\questiontitle{Sensitivity analysis vs.\ contribution analysis}
What is the main difference between sensitivity analysis (SA) and contribution analysis (CA)?
(Select one or more.)
\begin{itemize}[leftmargin=1.8em]
  \choice \textbf{SA} is gradient-based, while \textbf{CA} is not
  \choice \textbf{CA} focuses on how a small change in the input changes the output, and \textbf{SA} focuses on how the input is responsible for the output
  \choice \textbf{CA} is gradient-based, while \textbf{SA} is not
  \choice \textbf{SA} focuses on how a small change in the input changes the output, and \textbf{CA} focuses on how the input is responsible for the output
\end{itemize}
\closequestion


% ==========================================================
\section{Imitation Learning and Offline RL (mock)}
% ==========================================================

\questiontitle{Compounding errors in behavioral cloning}
A behavioral-cloning policy makes a mistake with probability $\epsilon$ under the expert's
state distribution. According to Ross \& Bagnell (2010), how does the total error grow with
the horizon $T$?
\begin{itemize}[leftmargin=1.8em]
  \choice Linearly, $\bigO{T\epsilon}$, because errors are independent across steps
  \choice Quadratically, $\bigO{T^2\epsilon}$, because an early mistake corrupts future states
  \choice It stays bounded by $\epsilon$ regardless of $T$
  \choice Logarithmically, $\bigO{\epsilon\log T}$
\end{itemize}
\closequestion

\questiontitle{What DAgger requires}
DAgger reduces the $\bigO{T^2}$ error of behavioral cloning to $\bigO{T}$. What is the central
cost or requirement that makes it applicable?
\begin{itemize}[leftmargin=1.8em]
  \choice A reward function for every state
  \choice An expert that can be \emph{queried for labels} on states the learner visits
  \choice A known transition model of the environment
  \choice A pretrained value function
\end{itemize}
\closequestion

\questiontitle{Multimodal demonstrations}
An expert demonstrates passing an obstacle on either the left or the right for similar states.
Why does a behavioral-cloning policy with an MSE loss fail here, and name one fix.
\workingspace{3.5cm}
\closequestion

\questiontitle{Why standard off-policy RL fails offline}
Running DQN or SAC on a fixed data set without new interaction typically diverges. What is the
mechanism?
\begin{itemize}[leftmargin=1.8em]
  \choice The replay buffer is too small to cover the state space
  \choice The $\max_{a'} Q(s',a')$ backup queries out-of-distribution actions, whose values the network overestimates, and the error propagates via bootstrapping
  \choice The learning rate must decay faster offline
  \choice Off-policy methods cannot use a discount factor offline
\end{itemize}
\closequestion

\questiontitle{Three families of offline RL}
Name the three families of offline-RL methods and the single principle they all enforce.
\workingspace{4cm}
\closequestion

\questiontitle{How IQL avoids OOD queries}
Implicit Q-Learning (IQL) replaces the $\max_{a'} Q(s',a')$ target with expectile regression
on a value function $V(s)$. Why does this avoid out-of-distribution action queries, and what
does $V(s)$ approximate as the expectile $\tau \to 1$?
\workingspace{3.5cm}
\closequestion

\questiontitle{Trajectory stitching}
Which offline-RL approaches can perform \emph{trajectory stitching} (combining sub-optimal
trajectories into a better-than-dataset policy)?
\begin{itemize}[leftmargin=1.8em]
  \choice The Decision Transformer, because it conditions on return-to-go
  \choice Bellman-backup methods (CQL, IQL), because value propagates through bootstrapping
  \choice Only behavior cloning
  \choice No offline method can stitch
\end{itemize}
\closequestion

\questiontitle{IL vs.\ offline RL}
State the key difference in objective between imitation learning and offline RL, and which one
can in principle exceed the quality of the data it is trained on.
\workingspace{3cm}
\closequestion


% ==========================================================
\section{Inverse Reinforcement Learning (real exam)}
% ==========================================================
\textit{\small From the Inverse RL exam, reproduced verbatim with original answer choices.}

\questiontitle{MaxEnt IRL with function approximation}
When using Maximum Entropy IRL with high-dimensional function approximation (e.g.\ neural
networks), one key bottleneck is:
\begin{itemize}[leftmargin=1.8em]
  \choice The inability to compute gradients with respect to $\psi$
  \choice The need to differentiate through a complex partition function or to run repeated forward--backward algorithms
  \choice The requirement that all expert trajectories be perfect
  \choice The immediate convergence toward a uniform distribution over all trajectories
\end{itemize}
\closequestion

\questiontitle{Why Maximum Entropy IRL?}
Why is Maximum Entropy IRL used?
\begin{itemize}[leftmargin=1.8em]
  \choice To ensure the reward function is constant across all states
  \choice Because the method aims to retain as much randomness as possible, subject to matching expert statistics
  \choice Because it cannot handle suboptimal demonstrations
  \choice Because a Dirac delta distribution over the best path is always infeasible
\end{itemize}
\closequestion

\questiontitle{A common challenge in IRL}
What is a common challenge when performing IRL?
\begin{itemize}[leftmargin=1.8em]
  \choice There is always a single unique reward function consistent with the expert's behavior
  \choice The learned reward is trivial to evaluate without further policy optimization
  \choice Expert demonstrations may be suboptimal, complicating reward inference
  \choice IRL does not require knowledge or assumptions about the environment's dynamics
\end{itemize}
\closequestion

\questiontitle{Unique aspect of IRL vs.\ imitation learning}
Compared to Imitation Learning (IL), a unique aspect of IRL is that IRL:
\begin{itemize}[leftmargin=1.8em]
  \choice Aims to learn the expert's reward rather than just copying actions
  \choice Only works in deterministic environments
  \choice Discards the idea of a reward function entirely
  \choice Ignores expert demonstrations
\end{itemize}
\closequestion

\questiontitle{Role of $p(\tau)$ in MaxEnt IRL}
In Maximum Entropy IRL, the probability of a trajectory $\tau$ is written as
$p_\psi(\tau \mid \mathcal{O}_{0:T}) \propto p(\tau)\exp(r_\psi(\tau))$. What is the role of
$p(\tau)$ in this expression?
\begin{itemize}[leftmargin=1.8em]
  \choice It enforces that all trajectories with the same reward get zero probability
  \choice It encodes the environment's dynamics or ``how likely'' the raw trajectory is, independent of reward
  \choice It ensures that the partition function is always $1$
  \choice It guarantees that expert demonstrations appear with higher probability than random trajectories by default
\end{itemize}
\closequestion

\questiontitle{Slack variable in maximum-margin IRL}
Which statement correctly describes a slack variable $\zeta$ in maximum-margin IRL?
\begin{itemize}[leftmargin=1.8em]
  \choice It controls how strictly the constraints must be satisfied, allowing the expert to be suboptimal
  \choice It ensures the maximum margin has infinite size
  \choice It forces the expert's trajectories to always have higher reward
  \choice It is irrelevant in IRL and only used in supervised SVMs
\end{itemize}
\closequestion

\questiontitle{What IRL starts with}
In IRL, we typically start with:
\begin{itemize}[leftmargin=1.8em]
  \choice A set of expert demonstrations (trajectories)
  \choice A fully known reward function from the environment
  \choice A labeled dataset with correct actions for each state
  \choice No knowledge about the environment's dynamics
\end{itemize}
\closequestion

\questiontitle{Interpreting the MaxEnt IRL gradient}
Consider the gradient of the Maximum Entropy IRL objective,
$\E_{\tau\sim\pi^*}[\grad_\psi r_\psi(\tau)] - \E_{\tau\sim p^\psi(\tau\mid\mathcal{O}_{0:T})}[\grad_\psi r_\psi(\tau)]$.
Conceptually, what does each term represent?
\begin{itemize}[leftmargin=1.8em]
  \choice The first term decreases the expert's reward, the second increases the current policy's reward
  \choice Both terms push the reward parameters in the same direction
  \choice The first term reinforces expert trajectories; the second penalizes the policy's own trajectories
  \choice They both approximate the state visitation frequencies
\end{itemize}
\closequestion

\questiontitle{What makes MaxEnt IRL log-likelihood hard}
When optimizing the log-likelihood in Maximum Entropy IRL, what makes this problem
computationally challenging?
\begin{itemize}[leftmargin=1.8em]
  \choice The partition function $Z$ involves a sum over all possible trajectories, which can be extremely large
  \choice The gradient of the likelihood does not depend on $\psi$
  \choice Each demonstration is guaranteed to be optimal
  \choice The objective always converges in one iteration
\end{itemize}
\closequestion

\questiontitle{Feature matching with a linear reward}
In feature-matching IRL with a linear reward $r_\psi(s,a) = \psi\T f(s,a)$, which statement is
true?
\begin{itemize}[leftmargin=1.8em]
  \choice The reward is assumed to be a non-linear function
  \choice We aim to match the expert's feature expectations $\mu(\pi)$
  \choice We never need to solve an MDP to update $\psi$
  \choice The reward depends only on the next state transition
\end{itemize}
\closequestion

\questiontitle{Purpose of the structured margin term}
The structured max-margin IRL formulation uses a loss term $D(\pi_{\text{expert}}, \pi)$ in
$\psi\T\mu(\pi_{\text{expert}}) > \max_\pi\bigl(\psi\T\mu(\pi) + D(\pi_{\text{expert}},\pi)\bigr)$.
What is the purpose of $D(\pi_{\text{expert}}, \pi)$?
\begin{itemize}[leftmargin=1.8em]
  \choice It ensures that only the final reward function matters
  \choice It differentiates between policies that are ``almost'' like the expert versus those that are very different
  \choice It nullifies any margin gained by suboptimal policies
  \choice It speeds up the solver by removing constraints
\end{itemize}
\closequestion

\questiontitle{Primary objective of max-margin IRL}
In the maximum-margin approach to IRL (inspired by SVMs), what is the primary objective?
\begin{itemize}[leftmargin=1.8em]
  \choice To find the widest margin in the raw observations of states
  \choice To maximize the margin between the expert's feature expectations and those of any other policy
  \choice To minimize the number of slack variables to zero
  \choice To guarantee that the expert policy is identical to all alternative policies
\end{itemize}
\closequestion

\questiontitle{Primary goal of IRL}
Which of the following is a primary goal of Inverse Reinforcement Learning?
\begin{itemize}[leftmargin=1.8em]
  \choice Recover the reward function from expert demonstrations
  \choice Perform clustering on reward signals
  \choice Directly mimic the expert's actions without any reward model
  \choice Randomize the learned policy to increase exploration
\end{itemize}
\closequestion


% ==========================================================
\section{Preference-Based RL, RLHF and DPO (mock)}
% ==========================================================

\questiontitle{The Bradley--Terry preference model}
A reward model is trained from pairwise preferences $\tau_w \succ \tau_l$. Write the
Bradley--Terry probability that $\tau_w$ is preferred in terms of the trajectory rewards
$r_\theta(\tau_w), r_\theta(\tau_l)$, and state what loss is used to fit it.
\workingspace{3cm}
\closequestion

\questiontitle{The KL penalty in RLHF}
The RLHF objective is $\max_\theta \E_{y\sim\pi_\theta}[\,r_\phi(x,y) - \beta\log\frac{\pi_\theta(y\mid x)}{\pi_{\text{ref}}(y\mid x)}\,]$.
What is the role of the KL term, and what failure mode does it mitigate?
\begin{itemize}[leftmargin=1.8em]
  \choice It increases the entropy of the reward model
  \choice It keeps the policy close to the reference (SFT) model, mitigating reward overoptimization
  \choice It guarantees the reward model is unbiased
  \choice It replaces the need for a value function in PPO
\end{itemize}
\closequestion

\questiontitle{Two anchors in PPO-style RLHF}
PPO-style RLHF uses two distinct reference points. Identify each and what it controls.
\workingspace{3.5cm}
\closequestion

\questiontitle{What DPO eliminates}
Direct Preference Optimization (DPO) rewrites the RLHF problem so that the reward and partition
function cancel. What does this eliminate compared to standard RLHF, and what is the main
limitation of DPO?
\begin{itemize}[leftmargin=1.8em]
  \choice It eliminates the SFT step; limitation: it needs more human labels
  \choice It eliminates the separate reward model and the RL loop, becoming supervised learning on preference pairs; limitation: it assumes a fixed offline preference set
  \choice It eliminates the KL penalty; limitation: it diverges
  \choice It eliminates pretraining; limitation: it cannot scale
\end{itemize}
\closequestion

\questiontitle{Goodhart's law and reward hacking}
Give the statement of Goodhart's law as it applies to RL, and one concrete example of reward
hacking discussed in the lecture.
\workingspace{3cm}
\closequestion

\questiontitle{GRPO vs.\ PPO}
Group Relative Policy Optimization (GRPO) avoids one expensive component of PPO-style RLHF.
Which component, and how does GRPO obtain its advantage estimate instead?
\begin{itemize}[leftmargin=1.8em]
  \choice It avoids the reward model; it uses human labels directly
  \choice It avoids the separate value (critic) network; it standardizes rewards within a group of sampled outputs for the same prompt
  \choice It avoids the KL penalty; it clips more aggressively
  \choice It avoids sampling; it uses a fixed dataset
\end{itemize}
\closequestion

\questiontitle{RLAIF and Constitutional AI}
What problem does RLAIF address, and what is the key insight that makes it work?
\workingspace{3cm}
\closequestion


% ==========================================================
\section{Planning and Search (mock)}
% ==========================================================

\questiontitle{The four phases of MCTS}
List the four phases of Monte Carlo Tree Search in order, and state which action is finally
played at the root.
\workingspace{4cm}
\closequestion

\questiontitle{The UCT tree policy}
UCT selects actions by $a^* = \argmax_a[\,Q(s,a) + c\sqrt{\ln N(s)/N(s,a)}\,]$. Explain what
each of the two terms encourages and the role of $c$.
\workingspace{3.5cm}
\closequestion

\questiontitle{AlphaZero's network}
How does AlphaZero differ from AlphaGo, and what does its single network $f_\theta(s)$ output?
\begin{itemize}[leftmargin=1.8em]
  \choice AlphaZero adds a human-games dataset; the network outputs only a value
  \choice AlphaZero removes human games, the rollout policy, and supervised pretraining; one network outputs both a policy prior $\mathbf{p}$ and a value $v$
  \choice AlphaZero replaces MCTS with $Q$-learning; the network outputs $Q$-values
  \choice AlphaZero uses two separate networks trained on expert data
\end{itemize}
\closequestion

\questiontitle{MCTS as policy improvement}
In AlphaZero, why is MCTS described as a ``policy-improvement operator,'' and what role does
the network play with respect to the search?
\workingspace{3.5cm}
\closequestion

\questiontitle{The branching problem}
At depth $d$ with branching factor $b$, a full search tree has how many nodes, and name the two
classical remedies MCTS uses to make the search tractable.
\answerline{4cm}
\closequestion

\questiontitle{Random shooting vs.\ CEM}
For continuous-control trajectory optimisation, contrast random shooting with the Cross-Entropy
Method (CEM). What does CEM do that random shooting does not?
\workingspace{3.5cm}
\closequestion

\questiontitle{Model Predictive Control}
In MPC, the agent optimizes an $H$-step action sequence but executes only part of it. Which
part, and why does this help when the model is imperfect?
\begin{itemize}[leftmargin=1.8em]
  \choice It executes the whole sequence open-loop; this is faster
  \choice It executes only the first action, then re-plans, which corrects for model errors after each real transition
  \choice It executes the last action; this maximizes lookahead
  \choice It executes a random action from the sequence for exploration
\end{itemize}
\closequestion


% ==========================================================
\section{Model-Based RL (mock)}
% ==========================================================

\questiontitle{Dyna-Q and the planning steps}
In Dyna-Q, real experience updates the model and the $Q$-function, while imagined experience
updates only the $Q$-function. What does the algorithm reduce to when the number of planning
steps $K = 0$?
\begin{itemize}[leftmargin=1.8em]
  \choice Monte Carlo control
  \choice Ordinary model-free $Q$-learning
  \choice Policy iteration
  \choice Behavioral cloning
\end{itemize}
\closequestion

\questiontitle{Why MBPO uses short rollouts}
MBPO starts model rollouts from real replay-buffer states and keeps them very short (often
$k=1$--$5$). Explain the trade-off that motivates short rollouts.
\workingspace{3.5cm}
\closequestion

\questiontitle{Aleatoric vs.\ epistemic uncertainty}
Define aleatoric and epistemic uncertainty, and state which one can be reduced by collecting
more data and which one matters most for model-based RL planning.
\workingspace{4cm}
\closequestion

\questiontitle{Ensembles and uncertainty}
A single Gaussian neural-network model outputs a mean and variance. Which type of uncertainty
does its predictive variance capture, and what must you do instead to estimate epistemic
uncertainty?
\begin{itemize}[leftmargin=1.8em]
  \choice It captures epistemic uncertainty; nothing more is needed
  \choice It captures aleatoric uncertainty; for epistemic uncertainty you need disagreement between an ensemble of models
  \choice It captures both equally; ensembles are unnecessary
  \choice It captures neither; only Bayesian neural networks work
\end{itemize}
\closequestion

\questiontitle{Model exploitation}
What is model exploitation in model-based RL, and name two mitigations.
\workingspace{3.5cm}
\closequestion

\questiontitle{PETS}
PETS combines three ingredients for continuous control. Name them and state what ``trajectory
sampling'' (rolling out particles through the ensemble) achieves.
\workingspace{3.5cm}
\closequestion

\questiontitle{Latent world models}
Why do agents like Dreamer and MuZero learn a \emph{latent} state $z_t$ instead of predicting
in the raw observation space? Give two reasons.
\workingspace{3.5cm}
\closequestion


% ==========================================================
\section{RL for Foundation Models (mock)}
% ==========================================================

\questiontitle{Autoregressive generation as RL}
Map an autoregressive language model onto the RL formalism: what plays the role of the policy,
the state, the action, and the reward?
\workingspace{3.5cm}
\closequestion

\questiontitle{Pretraining as imitation}
Why is language-model pretraining described as a form of behavioral cloning, and why is it
``not enough'' to produce a helpful assistant?
\workingspace{3.5cm}
\closequestion

\questiontitle{GRPO advantage normalization}
In GRPO, $G$ outputs are sampled for the same prompt $x$ and the advantage is
$\hat A_i = (r_i - \mathrm{mean}(r_{1:G}))/(\mathrm{std}(r_{1:G}) + \delta)$. At what level
(token or sequence) is $\hat A_i$ computed, and how is it applied to the tokens of output $y_i$?
\begin{itemize}[leftmargin=1.8em]
  \choice Token-level; applied independently to each token
  \choice Sequence-level; the same $\hat A_i$ is applied to all tokens of output $y_i$
  \choice It is not an advantage at all; it is a reward
  \choice It is computed across prompts, not within a group
\end{itemize}
\closequestion

\questiontitle{Verifiable rewards (RLVR)}
What is the key idea of Reinforcement Learning with Verifiable Rewards (RLVR), and why is a
composite reward $r(y) = \lambda_{\text{format}} r_{\text{format}}(y) + \lambda_{\text{correct}} r_{\text{correct}}(y)$
often used?
\workingspace{4cm}
\closequestion

\questiontitle{SFT vs.\ RLHF vs.\ RLVR}
For each of SFT, RLHF, and RLVR, give the training signal it relies on and its main risk.
\workingspace{4.5cm}
\closequestion

\questiontitle{Generalist agents}
Gato uses one Transformer with the same weights across Atari, robotics, captioning, and
dialogue. What makes this possible, and how was Gato primarily trained?
\workingspace{3.5cm}
\closequestion


% ==========================================================
%  ANSWER KEY
% ==========================================================
\newpage
\section*{Answer Key}
\addcontentsline{toc}{section}{Answer Key}

\textit{\small Real-exam answers (Q1--12 credit assignment, Q21--33 IRL) follow the official keys;
mock answers give the intended solution.}

\subsection*{Credit Assignment (real)}
\begin{enumerate}[label=\textbf{Q\arabic*.}, leftmargin=3em, itemsep=3pt]
  \item \textbf{C} --- a strictly monotonic transformation of the \emph{total return} preserves the argmax over policies; affine/clipping/squaring of per-step rewards need not.
  \item \textbf{D} --- potential-based shaping with $\Phi$ is equivalent to initializing $Q$-values with the potential function.
  \item \textbf{A} --- the LSTM's consecutive prediction errors are highly correlated, so they cancel in the difference, giving lower variance.
  \item \textbf{A} --- if $\Phi$ is learned online from sparse, delayed TD-errors, it must itself be bootstrapped, which re-introduces bias.
  \item \textbf{C} --- credit assignment is assigning credit/blame for overall outcomes to each of the decisions that contributed.
  \item \textbf{A, C} --- variance in MC grows with reward delay; TD($\lambda$) assigns credit faster than TD($0$).
  \item \textbf{C} --- with zero expected future reward the reward depends on the previous state/action, so it is no longer Markov.
  \item \textbf{C} --- an SDP has a Markov policy and Markov transitions, but the reward need not be Markov.
  \item \textbf{B} --- return-equivalent SDPs differ only in reward distribution but have the same \emph{expected return}.
  \item \textbf{A} --- because later states are used for prediction, earlier reward-causing actions get explained away (missed).
  \item \textbf{A, B, C, D, E} --- all statements about LSTM vs.\ FFN prediction are correct.
  \item \textbf{D} --- SA measures how a small input change affects the output; CA measures how the input is \emph{responsible} for the output.
\end{enumerate}

\subsection*{Imitation Learning and Offline RL (mock)}
\begin{enumerate}[label=\textbf{Q\arabic*.}, leftmargin=3em, itemsep=3pt, start=13]
  \item \textbf{B} --- error grows quadratically, $\bigO{T^2\epsilon}$, because mistakes change future states.
  \item \textbf{B} --- DAgger needs an expert that can label the learner's on-policy states.
  \item MSE \emph{averages} the two valid modes ($\approx\frac12(a_{\text{left}}+a_{\text{right}})$), driving into the obstacle. Fix: mixture-of-Gaussians head, latent-variable model (CVAE/flow), or autoregressive action discretization.
  \item \textbf{B} --- the $\max_{a'}Q$ backup queries OOD actions; the net overestimates them and the error propagates through bootstrapping.
  \item Behavior regularization (stay close to $\pi_\beta$), conservative value estimation (CQL), and sequence modeling (Decision Transformer). Common principle: \emph{do not trust the model at out-of-distribution inputs.}
  \item IQL regresses $V(s)$ using \emph{expectile} loss over \emph{dataset actions only}, so no unseen action is ever evaluated; as $\tau\to1$, $V(s)$ approaches a soft $\max_{a\in\dataset}Q(s,a)$.
  \item \textbf{B} --- Bellman-backup methods (CQL, IQL) can stitch; the Decision Transformer cannot.
  \item IL matches the expert (upper bound: expert level); offline RL maximizes reward and can in principle exceed the dataset's behavior policy.
\end{enumerate}

\subsection*{Inverse Reinforcement Learning (real)}
\begin{enumerate}[label=\textbf{Q\arabic*.}, leftmargin=3em, itemsep=3pt, start=21]
  \item \textbf{B} --- differentiating through the partition function / repeated forward--backward passes.
  \item \textbf{B} --- retain maximum randomness subject to matching expert statistics.
  \item \textbf{C} --- expert demonstrations may be suboptimal.
  \item \textbf{A} --- IRL learns the expert's reward rather than copying actions.
  \item \textbf{B} --- $p(\tau)$ encodes environment dynamics / how likely the raw trajectory is, independent of reward.
  \item \textbf{A} --- the slack variable relaxes the constraints, allowing a suboptimal expert.
  \item \textbf{A} --- IRL starts with a set of expert demonstrations (trajectories).
  \item \textbf{C} --- first term reinforces expert trajectories; second penalizes the policy's own trajectories.
  \item \textbf{A} --- the partition function $Z$ sums over all possible trajectories.
  \item \textbf{B} --- we aim to match the expert's feature expectations $\mu(\pi)$.
  \item \textbf{B} --- $D$ differentiates ``almost-expert'' policies from very different ones (structured margin).
  \item \textbf{B} --- maximize the margin between the expert's feature expectations and those of any other policy.
  \item \textbf{A} --- recover the reward function from expert demonstrations.
\end{enumerate}

\subsection*{Preference-Based RL, RLHF and DPO (mock)}
\begin{enumerate}[label=\textbf{Q\arabic*.}, leftmargin=3em, itemsep=3pt, start=34]
  \item $P(\tau_w\succ\tau_l)=\sigma(r_\theta(\tau_w)-r_\theta(\tau_l))$, with $r_\theta(\tau)=\sum_{(s,a)\in\tau}r_\theta(s,a)$; fit by maximizing log-likelihood (binary cross-entropy on the preference label).
  \item \textbf{B} --- the KL term keeps $\pi_\theta$ near the SFT reference, mitigating reward overoptimization.
  \item KL penalty: $\pi_\theta$ vs.\ $\pi_{\text{ref}}$ (SFT) --- stay close to the trusted assistant. PPO clipping: $\pi_\theta$ vs.\ $\pi_{\theta_{\text{old}}}$ --- keep each update step stable.
  \item \textbf{B} --- DPO removes the separate reward model and the RL loop (pure supervised learning on pairs); limitation: it assumes a fixed offline preference set.
  \item ``When a measure becomes a target, it ceases to be a good measure.'' Example: the OpenAI boat-race agent spinning to collect speed boosts instead of finishing (or a verbose/sycophantic RLHF model; a robot moving the camera so an object only \emph{appears} grasped).
  \item \textbf{B} --- GRPO drops the value/critic network and standardizes rewards within a group of sampled outputs for the same prompt.
  \item RLAIF replaces human annotators with a strong LLM because human labelling does not scale; key insight: critique is easier than generation. Constitutional AI encodes the values as written principles the model uses to critique and revise its own outputs.
\end{enumerate}

\subsection*{Planning and Search (mock)}
\begin{enumerate}[label=\textbf{Q\arabic*.}, leftmargin=3em, itemsep=3pt, start=41]
  \item Selection (tree policy / UCT) $\to$ Expansion $\to$ Simulation (rollout) $\to$ Backpropagation. The action with the \emph{highest visit count} at the root is played.
  \item $Q(s,a)$ exploits promising branches; $c\sqrt{\ln N(s)/N(s,a)}$ explores under-visited actions (large when $N(s,a)$ is small). $c$ sets the exploration--exploitation balance.
  \item \textbf{B} --- AlphaZero removes human games, the rollout policy, and supervised pretraining; one network outputs both a policy prior $\mathbf{p}$ and a value $v$.
  \item MCTS turns the raw network policy $\mathbf{p}_\theta$ into a stronger search policy $\pi_{\text{MCTS}}\propto N(s,a)^{1/\tau}$; the network is then trained to \emph{amortize} (imitate) that improved search and the final outcome.
  \item $b^d$ nodes. Remedies: depth reduction (a learned value function replaces deep rollouts) and breadth reduction (expand only promising branches).
  \item Random shooting samples action sequences uniformly and keeps the best; CEM \emph{iteratively refits} the sampling distribution to the elite (best-scoring) sequences, concentrating samples on good regions.
  \item \textbf{B} --- MPC executes only the first action, then re-plans, correcting model errors after each real transition.
\end{enumerate}

\subsection*{Model-Based RL (mock)}
\begin{enumerate}[label=\textbf{Q\arabic*.}, leftmargin=3em, itemsep=3pt, start=48]
  \item \textbf{B} --- with $K=0$ Dyna-Q is ordinary model-free $Q$-learning.
  \item Short rollouts trade off model usefulness against compounding error: longer rollouts give more imagined data but accumulate model error and invite exploitation. Starting from real states and imagining only a few steps keeps synthetic data trustworthy. Goal: use the model only where it is accurate.
  \item Aleatoric = inherent environment randomness (not reducible with more data); epistemic = uncertainty from limited data / OOD queries (reducible). Epistemic uncertainty matters most for MBRL --- it tells the planner where the model should not be trusted.
  \item \textbf{B} --- a single Gaussian net captures aleatoric uncertainty (predictive variance); epistemic uncertainty requires disagreement across an ensemble.
  \item Model exploitation: a policy optimized against a learned model finds and exploits its errors (e.g.\ a glitch giving infinite simulated reward). Mitigations: limit the planning horizon, penalize uncertainty (pessimistic planning), and frequently retrain on new real data.
  \item Probabilistic ensemble (dynamics + uncertainty), CEM (action-sequence search), and trajectory sampling (ensemble rollouts). Trajectory sampling propagates model disagreement into the planning, so uncertain actions are down-weighted.
  \item Raw observations are high-dimensional, partial, contain irrelevant variation, and may miss history-dependent information; a latent $z_t$ summarizes the useful information for control (the goal is a control-useful representation, not perfect reconstruction).
\end{enumerate}

\subsection*{RL for Foundation Models (mock)}
\begin{enumerate}[label=\textbf{Q\arabic*.}, leftmargin=3em, itemsep=3pt, start=54]
  \item Policy = $\pi_\theta(y_t\mid x,y_{<t})$ (the LM); state = prompt + generated prefix; action = next token (or tool call); reward = preference / verifier / task-success signal.
  \item Pretraining maximizes token likelihood on text, i.e.\ it \emph{imitates} the conditional distribution of the data --- behavioral cloning. It is not enough because it learns plausible (and conflicting) continuations, not instruction-following, calibration, refusal, or alignment with preferences.
  \item \textbf{B} --- $\hat A_i$ is sequence-level (one value per output) and the same $\hat A_i$ is applied to all tokens of $y_i$.
  \item RLVR replaces a learned preference reward with an automatic checker (math verifier, unit tests, API check, game score). The composite reward teaches both \emph{format} (answer must be extractable/parseable) and \emph{correctness} (extracted answer passes the verifier); without a format term, capability is underestimated.
  \item SFT: demonstrations --- risk: limited by examples. RLHF: human preferences --- risk: proxy reward hacking / overoptimization. RLVR: automatic verifier --- risk: narrow rewards and hacking the verifier.
  \item Tokenizing all modalities (text, vision, state, actions) into a shared sequence makes language modeling and action prediction structurally identical, so one Transformer can serve all tasks. Gato was trained primarily by supervised sequence modeling (imitation) on mixed datasets.
\end{enumerate}


\end{document}
